\documentclass[11pt,a4paper]{article}
\usepackage[utf8]{inputenc}
\usepackage[T1]{fontenc}
\usepackage[spanish,es-minimal,es-noshorthands]{babel}
\usepackage{lmodern}
\usepackage[sfdefault,lining]{FiraSans}
\usepackage[a4paper,top=2.3cm,bottom=2.2cm,left=1.8cm,right=1.8cm,headheight=1.3cm,headsep=0.4cm,footskip=1cm]{geometry}
\usepackage{amsmath,amssymb}
\usepackage{icomma}
\usepackage{graphicx}
\usepackage[table]{xcolor}
\usepackage{tikz}
\usepackage{fontawesome5}
\usepackage{ccicons}
\usepackage{fancyhdr}
\usepackage{lastpage}
\usepackage{tasks}
\usepackage{needspace}
\usepackage[most]{tcolorbox}
\usepackage{array}
\usepackage[hidelinks]{hyperref}
\AtBeginDocument{\hypersetup{pdftitle={Ficha de ejercicios de la unidad 4: Geometría analítica en el plano. 1º Bach.},pdfsubject={Matemáticas I 1º Bach.},pdfkeywords={geometría analítica, vectores, producto escalar, ecuaciones de la recta, posiciones relativas, rectas perpendiculares, ángulos, distancias, mediatriz, bisectriz}}}

\definecolor{navy}{HTML}{1B2748}
\definecolor{accent}{HTML}{E8603A}
\definecolor{teal}{HTML}{1F9E90}
\definecolor{soft}{HTML}{EEF1F8}
\definecolor{muted}{HTML}{5D6685}

\setlength{\parindent}{0pt}
\setlength{\parskip}{3pt}
\settasks{label=\alph*), label-format=\bfseries\color{navy}, label-width=1.3em, item-indent=1.8em, column-sep=0.6em, after-item-skip=5pt, before-skip=3pt, after-skip=4pt}

% ---------- capçalera i peu ----------
\newcommand{\llicencia}{\href{https://creativecommons.org/licenses/by-nc-sa/4.0/deed.es}{\ccbyncsa}\ \ Licencia Creative Commons Atribución-NoComercial-CompartirIgual 4.0 Internacional (CC BY-NC-SA 4.0)}
\pagestyle{fancy}
\fancyhf{}
\fancyhead[L]{\raisebox{-0.35\height}{\includegraphics[height=0.95cm]{logo-gva.png}}}
\fancyhead[C]{\small\color{muted}Ficha de ejercicios · Unidad 4: Geometría analítica en el plano · 1\textsuperscript{o} Bach.}
\fancyhead[R]{\raisebox{-0.35\height}{\includegraphics[height=0.95cm]{logo-vermellar.png}}}
\fancyfoot[L]{\scriptsize\color{muted}\llicencia}
\fancyfoot[R]{\small\color{muted}\thepage}
\renewcommand{\headrulewidth}{0.4pt}
\renewcommand{\headrule}{\color{muted!40}\hrule height \headrulewidth}
\fancypagestyle{primera}{\fancyhf{}\renewcommand{\headrulewidth}{0pt}
  \fancyfoot[L]{\scriptsize\color{muted}\llicencia}
  \fancyfoot[R]{\small\color{muted}\thepage}}

% ---------- nivells i exercicis ----------
\newcommand{\estrelles}[1]{{\color{accent}\ifcase#1\or\faStar\faStar[regular]\faStar[regular]\or\faStar\faStar\faStar[regular]\or\faStar\faStar\faStar\fi}}
\newcounter{exer}
\newcommand{\ex}[1]{\par\needspace{4\baselineskip}\medskip\refstepcounter{exer}\noindent{\color{navy}\bfseries\theexer.}\ \ }

\newcommand{\seccio}[2]{\par\needspace{16\baselineskip}\bigskip
  \begin{tcolorbox}[colback=navy,colframe=navy,coltext=white,arc=2mm,boxrule=0pt,left=3mm,right=3mm,top=1.2mm,bottom=1.2mm]
  {\large\bfseries #1\quad #2}
  \end{tcolorbox}\vspace{-2pt}}

\newcommand{\exf}[1]{\par\medskip\refstepcounter{exer}\noindent{\color{navy}\bfseries\theexer.}\ \ }

% quadre de recordatori (només a la fitxa d'exercicis)
\newcommand{\recorda}[1]{\par\smallskip\begin{tcolorbox}[colback=accent!7,colframe=accent!55,arc=2mm,boxrule=0.5pt,left=2.5mm,right=2.5mm,top=1mm,bottom=1mm]\small{\color{accent}\faLightbulb[regular]}\ \textbf{Recuerda:}\ #1\end{tcolorbox}}
\newcommand{\res}[1]{#1}\newcommand{\rest}[1]{#1}

\def\fmneed{11}% que el títol d'una secció no quede sol al final del full
% =====================================================================
% Capçalera FentMates de les fitxes: un color per a cada curs.
% Els colors són els de la pàgina inicial (LEVELS a assets/js/materials.js);
% el logo ve de _eines/logo/pdf/ (python3 _eines/logo/genera.py).
%
% Ús, al final del preàmbul de la unitat:
%   % =====================================================================
% Capçalera FentMates de les fitxes: un color per a cada curs.
% Els colors són els de la pàgina inicial (LEVELS a assets/js/materials.js);
% el logo ve de _eines/logo/pdf/ (python3 _eines/logo/genera.py).
%
% Ús, al final del preàmbul de la unitat:
%   % =====================================================================
% Capçalera FentMates de les fitxes: un color per a cada curs.
% Els colors són els de la pàgina inicial (LEVELS a assets/js/materials.js);
% el logo ve de _eines/logo/pdf/ (python3 _eines/logo/genera.py).
%
% Ús, al final del preàmbul de la unitat:
%   \input{../comu/fentmates.tex}
%   \fmcurs{1eso}
% i, al començament del document:
%   \capcalera{Matemàtiques}{1r ESO}{Fitxa d'exercicis · Unitat 2}{Divisibilitat}
% Opcional abans de \input: \def\fmneed{14} (espai mínim abans de cada secció).
% =====================================================================
\makeatletter
\def\fm@logodir{../../logo/pdf}
\providecommand{\fmneed}{10}
\def\fm@cursnom{}

% clau del curs, color i color del text damunt del color
\newcommand{\fm@curs}[3]{\@namedef{fm@color@#1}{#2}\@namedef{fm@text@#1}{#3}}
\fm@curs{1eso}{E6B84A}{navy}
\fm@curs{2eso}{3B82F6}{white}
\fm@curs{3eso}{FF7A50}{white}
\fm@curs{4eso-a}{A78BFA}{white}
\fm@curs{4eso-b}{6D28D9}{white}
\fm@curs{1batx-mi}{9B1C1C}{white}
\fm@curs{1batx-ccss}{EF6B6B}{white}
\fm@curs{2batx-mii}{1E6B47}{white}
\fm@curs{2batx-ccss}{2BB3A3}{white}

\newcommand{\fmcurs}[1]{%
  \@ifundefined{fm@color@#1}{\PackageError{fentmates}{Curs desconegut: #1}{}}{}%
  \edef\fm@tmp{\noexpand\definecolor{curs}{HTML}{\@nameuse{fm@color@#1}}%
               \noexpand\colorlet{curstext}{\@nameuse{fm@text@#1}}}\fm@tmp
  \colorlet{cursfosc}{curs!72!black}%
  \def\fm@logo{\fm@logodir/vertical-#1.pdf}%
  \def\fm@marca{\fm@logodir/marca-#1.pdf}}

% ---------- capçalera de la primera pàgina ----------
% #1 matèria (petit), #2 curs, #3 tipus de document i unitat, #4 títol
\newcommand{\capcalera}[4]{%
  \gdef\fm@cursnom{#2}%
  \par\noindent
  \begin{tcolorbox}[colback=curs,colframe=curs,boxrule=0pt,arc=2.2mm,width=3.5cm,height=2.15cm,
      valign=center,left=3mm,right=2mm,top=1mm,bottom=1mm,nobeforeafter,box align=center]
    \color{curstext}{\footnotesize #1}\par\vspace{1pt}{\fontsize{19}{21}\selectfont\bfseries #2\par}%
  \end{tcolorbox}\hspace{5mm}%
  \begin{minipage}[c]{\dimexpr\linewidth-3.5cm-5mm-2.9cm\relax}
    {\color{cursfosc}\small\bfseries\MakeUppercase{#3}}\par\vspace{2pt}
    {\color{navy}\fontsize{20}{23}\selectfont\bfseries #4\par}
  \end{minipage}\hfill
  \begin{minipage}[c]{2.6cm}\raggedleft\includegraphics[width=2.6cm]{\fm@logo}\end{minipage}\par
  \vspace{2pt}}

% ---------- barres de secció del color del curs ----------
\renewcommand{\seccio}[2]{\par\needspace{\fmneed\baselineskip}\bigskip
  \begin{tcolorbox}[colback=curs,colframe=curs,coltext=curstext,arc=2mm,boxrule=0pt,left=3mm,right=3mm,top=1.2mm,bottom=1.2mm]
  {\large\bfseries #1\quad #2}
  \end{tcolorbox}\vspace{-2pt}}

% ---------- estrelles de dificultat del color del curs (a Batxillerat no se'n posen) ----------
\@ifundefined{estrelles}{}{%
  \renewcommand{\estrelles}[1]{{\color{curs}\ifcase#1\or\faStar\faStar[regular]\faStar[regular]\or\faStar\faStar\faStar[regular]\or\faStar\faStar\faStar\fi}}}

% ---------- capçalera i peu de les pàgines ----------
% primera pàgina: logos institucionals al peu
\fancypagestyle{primera}{\fancyhf{}\renewcommand{\headrulewidth}{0pt}
  \fancyfoot[L]{\raisebox{-0.32\height}{\includegraphics[height=0.78cm]{logo-gva.png}}\hspace{5mm}%
    \raisebox{-0.32\height}{\includegraphics[height=0.78cm]{logo-vermellar.png}}}
  \fancyfoot[R]{\scriptsize\color{muted}\href{https://creativecommons.org/licenses/by-nc-sa/4.0/deed.ca}{\ccbyncsa}\ \ CC BY-NC-SA 4.0 · fentmates.org\qquad\small\thepage}}
% la resta: el símbol a l'esquerra i el curs a la dreta, amb el títol al centre com abans
\fancyhead[L]{\raisebox{-0.3\height}{\includegraphics[height=0.62cm]{\fm@marca}}}
\fancyhead[R]{\small\bfseries\color{cursfosc}\fm@cursnom}
\renewcommand{\headrule}{\color{curs!45}\hrule height \headrulewidth}
\makeatother

%   \fmcurs{1eso}
% i, al començament del document:
%   \capcalera{Matemàtiques}{1r ESO}{Fitxa d'exercicis · Unitat 2}{Divisibilitat}
% Opcional abans de \input: \def\fmneed{14} (espai mínim abans de cada secció).
% =====================================================================
\makeatletter
\def\fm@logodir{../../logo/pdf}
\providecommand{\fmneed}{10}
\def\fm@cursnom{}

% clau del curs, color i color del text damunt del color
\newcommand{\fm@curs}[3]{\@namedef{fm@color@#1}{#2}\@namedef{fm@text@#1}{#3}}
\fm@curs{1eso}{E6B84A}{navy}
\fm@curs{2eso}{3B82F6}{white}
\fm@curs{3eso}{FF7A50}{white}
\fm@curs{4eso-a}{A78BFA}{white}
\fm@curs{4eso-b}{6D28D9}{white}
\fm@curs{1batx-mi}{9B1C1C}{white}
\fm@curs{1batx-ccss}{EF6B6B}{white}
\fm@curs{2batx-mii}{1E6B47}{white}
\fm@curs{2batx-ccss}{2BB3A3}{white}

\newcommand{\fmcurs}[1]{%
  \@ifundefined{fm@color@#1}{\PackageError{fentmates}{Curs desconegut: #1}{}}{}%
  \edef\fm@tmp{\noexpand\definecolor{curs}{HTML}{\@nameuse{fm@color@#1}}%
               \noexpand\colorlet{curstext}{\@nameuse{fm@text@#1}}}\fm@tmp
  \colorlet{cursfosc}{curs!72!black}%
  \def\fm@logo{\fm@logodir/vertical-#1.pdf}%
  \def\fm@marca{\fm@logodir/marca-#1.pdf}}

% ---------- capçalera de la primera pàgina ----------
% #1 matèria (petit), #2 curs, #3 tipus de document i unitat, #4 títol
\newcommand{\capcalera}[4]{%
  \gdef\fm@cursnom{#2}%
  \par\noindent
  \begin{tcolorbox}[colback=curs,colframe=curs,boxrule=0pt,arc=2.2mm,width=3.5cm,height=2.15cm,
      valign=center,left=3mm,right=2mm,top=1mm,bottom=1mm,nobeforeafter,box align=center]
    \color{curstext}{\footnotesize #1}\par\vspace{1pt}{\fontsize{19}{21}\selectfont\bfseries #2\par}%
  \end{tcolorbox}\hspace{5mm}%
  \begin{minipage}[c]{\dimexpr\linewidth-3.5cm-5mm-2.9cm\relax}
    {\color{cursfosc}\small\bfseries\MakeUppercase{#3}}\par\vspace{2pt}
    {\color{navy}\fontsize{20}{23}\selectfont\bfseries #4\par}
  \end{minipage}\hfill
  \begin{minipage}[c]{2.6cm}\raggedleft\includegraphics[width=2.6cm]{\fm@logo}\end{minipage}\par
  \vspace{2pt}}

% ---------- barres de secció del color del curs ----------
\renewcommand{\seccio}[2]{\par\needspace{\fmneed\baselineskip}\bigskip
  \begin{tcolorbox}[colback=curs,colframe=curs,coltext=curstext,arc=2mm,boxrule=0pt,left=3mm,right=3mm,top=1.2mm,bottom=1.2mm]
  {\large\bfseries #1\quad #2}
  \end{tcolorbox}\vspace{-2pt}}

% ---------- estrelles de dificultat del color del curs (a Batxillerat no se'n posen) ----------
\@ifundefined{estrelles}{}{%
  \renewcommand{\estrelles}[1]{{\color{curs}\ifcase#1\or\faStar\faStar[regular]\faStar[regular]\or\faStar\faStar\faStar[regular]\or\faStar\faStar\faStar\fi}}}

% ---------- capçalera i peu de les pàgines ----------
% primera pàgina: logos institucionals al peu
\fancypagestyle{primera}{\fancyhf{}\renewcommand{\headrulewidth}{0pt}
  \fancyfoot[L]{\raisebox{-0.32\height}{\includegraphics[height=0.78cm]{logo-gva.png}}\hspace{5mm}%
    \raisebox{-0.32\height}{\includegraphics[height=0.78cm]{logo-vermellar.png}}}
  \fancyfoot[R]{\scriptsize\color{muted}\href{https://creativecommons.org/licenses/by-nc-sa/4.0/deed.ca}{\ccbyncsa}\ \ CC BY-NC-SA 4.0 · fentmates.org\qquad\small\thepage}}
% la resta: el símbol a l'esquerra i el curs a la dreta, amb el títol al centre com abans
\fancyhead[L]{\raisebox{-0.3\height}{\includegraphics[height=0.62cm]{\fm@marca}}}
\fancyhead[R]{\small\bfseries\color{cursfosc}\fm@cursnom}
\renewcommand{\headrule}{\color{curs!45}\hrule height \headrulewidth}
\makeatother

%   \fmcurs{1eso}
% i, al començament del document:
%   \capcalera{Matemàtiques}{1r ESO}{Fitxa d'exercicis · Unitat 2}{Divisibilitat}
% Opcional abans de \input: \def\fmneed{14} (espai mínim abans de cada secció).
% =====================================================================
\makeatletter
\def\fm@logodir{../../logo/pdf}
\providecommand{\fmneed}{10}
\def\fm@cursnom{}

% clau del curs, color i color del text damunt del color
\newcommand{\fm@curs}[3]{\@namedef{fm@color@#1}{#2}\@namedef{fm@text@#1}{#3}}
\fm@curs{1eso}{E6B84A}{navy}
\fm@curs{2eso}{3B82F6}{white}
\fm@curs{3eso}{FF7A50}{white}
\fm@curs{4eso-a}{A78BFA}{white}
\fm@curs{4eso-b}{6D28D9}{white}
\fm@curs{1batx-mi}{9B1C1C}{white}
\fm@curs{1batx-ccss}{EF6B6B}{white}
\fm@curs{2batx-mii}{1E6B47}{white}
\fm@curs{2batx-ccss}{2BB3A3}{white}

\newcommand{\fmcurs}[1]{%
  \@ifundefined{fm@color@#1}{\PackageError{fentmates}{Curs desconegut: #1}{}}{}%
  \edef\fm@tmp{\noexpand\definecolor{curs}{HTML}{\@nameuse{fm@color@#1}}%
               \noexpand\colorlet{curstext}{\@nameuse{fm@text@#1}}}\fm@tmp
  \colorlet{cursfosc}{curs!72!black}%
  \def\fm@logo{\fm@logodir/vertical-#1.pdf}%
  \def\fm@marca{\fm@logodir/marca-#1.pdf}}

% ---------- capçalera de la primera pàgina ----------
% #1 matèria (petit), #2 curs, #3 tipus de document i unitat, #4 títol
\newcommand{\capcalera}[4]{%
  \gdef\fm@cursnom{#2}%
  \par\noindent
  \begin{tcolorbox}[colback=curs,colframe=curs,boxrule=0pt,arc=2.2mm,width=3.5cm,height=2.15cm,
      valign=center,left=3mm,right=2mm,top=1mm,bottom=1mm,nobeforeafter,box align=center]
    \color{curstext}{\footnotesize #1}\par\vspace{1pt}{\fontsize{19}{21}\selectfont\bfseries #2\par}%
  \end{tcolorbox}\hspace{5mm}%
  \begin{minipage}[c]{\dimexpr\linewidth-3.5cm-5mm-2.9cm\relax}
    {\color{cursfosc}\small\bfseries\MakeUppercase{#3}}\par\vspace{2pt}
    {\color{navy}\fontsize{20}{23}\selectfont\bfseries #4\par}
  \end{minipage}\hfill
  \begin{minipage}[c]{2.6cm}\raggedleft\includegraphics[width=2.6cm]{\fm@logo}\end{minipage}\par
  \vspace{2pt}}

% ---------- barres de secció del color del curs ----------
\renewcommand{\seccio}[2]{\par\needspace{\fmneed\baselineskip}\bigskip
  \begin{tcolorbox}[colback=curs,colframe=curs,coltext=curstext,arc=2mm,boxrule=0pt,left=3mm,right=3mm,top=1.2mm,bottom=1.2mm]
  {\large\bfseries #1\quad #2}
  \end{tcolorbox}\vspace{-2pt}}

% ---------- estrelles de dificultat del color del curs (a Batxillerat no se'n posen) ----------
\@ifundefined{estrelles}{}{%
  \renewcommand{\estrelles}[1]{{\color{curs}\ifcase#1\or\faStar\faStar[regular]\faStar[regular]\or\faStar\faStar\faStar[regular]\or\faStar\faStar\faStar\fi}}}

% ---------- capçalera i peu de les pàgines ----------
% primera pàgina: logos institucionals al peu
\fancypagestyle{primera}{\fancyhf{}\renewcommand{\headrulewidth}{0pt}
  \fancyfoot[L]{\raisebox{-0.32\height}{\includegraphics[height=0.78cm]{logo-gva.png}}\hspace{5mm}%
    \raisebox{-0.32\height}{\includegraphics[height=0.78cm]{logo-vermellar.png}}}
  \fancyfoot[R]{\scriptsize\color{muted}\href{https://creativecommons.org/licenses/by-nc-sa/4.0/deed.ca}{\ccbyncsa}\ \ CC BY-NC-SA 4.0 · fentmates.org\qquad\small\thepage}}
% la resta: el símbol a l'esquerra i el curs a la dreta, amb el títol al centre com abans
\fancyhead[L]{\raisebox{-0.3\height}{\includegraphics[height=0.62cm]{\fm@marca}}}
\fancyhead[R]{\small\bfseries\color{cursfosc}\fm@cursnom}
\renewcommand{\headrule}{\color{curs!45}\hrule height \headrulewidth}
\makeatother

\fmcurs{1batx-mi}

\begin{document}
\thispagestyle{primera}

\capcalera{Matemáticas I}{1\textsuperscript{o} Bach.}{Ficha de ejercicios · Unidad 4}{\raggedright Geometría analítica en el plano}
\vspace{10pt}
Nombre y apellidos: \rule{7.8cm}{0.4pt}\hfill Grupo: \rule{1.6cm}{0.4pt}\hfill Fecha: \rule{2.2cm}{0.4pt}
\vspace{8pt}

\begin{tcolorbox}[colback=soft,colframe=soft,arc=2mm,left=3mm,right=3mm,top=1.5mm,bottom=1.5mm]
\small Resuelve los ejercicios en el cuaderno mostrando todos los pasos. Haz un dibujo aproximado en unos ejes de coordenadas siempre que te ayude: te permite detectar errores. Da los resultados de manera exacta (fracciones y raíces) y, si hace falta, también aproximados. Las ecuaciones de las rectas, en forma general $Ax+By+C=0$ con coeficientes enteros, si no se dice lo contrario.
\end{tcolorbox}

% =====================================================
\seccio{1}{Vectores y operaciones con vectores}

\exf{1} Dados los puntos $A(3,\,-1)$ y $B(-2,\,4)$, calcula:\begin{tasks}(2)
  \task Las componentes de $\overrightarrow{AB}$ y de $\overrightarrow{BA}$.
  \task El módulo de $\overrightarrow{AB}$.
  \task Un vector unitario con la misma dirección y sentido que $\overrightarrow{AB}$.
  \task Un vector de módulo 10 con la misma dirección y sentido contrario.
\end{tasks}

\ex{1} Calcula el argumento (el ángulo con el semieje positivo $OX$) de cada vector. ¡Fíjate en el cuadrante!\begin{tasks}(3)
  \task $\vec a=(3,\,3)$
  \task $\vec a=(-1,\,1)$
  \task $\vec a=(-4,\,-4)$
  \task $\vec a=(0,\,-3)$
  \task $\vec a=(5,\,0)$
  \task $\vec a=(2,\,-2)$
\end{tasks}

\ex{1} Calcula las componentes de un vector $\vec v$ si:\begin{tasks}(3)
  \task $|\vec v|=4$, argumento $30^\circ$
  \task $|\vec v|=10$, argumento $135^\circ$
  \task $|\vec v|=6$, argumento $240^\circ$
\end{tasks}\recorda{$\vec v=\big(|\vec v|\cos\alpha,\ |\vec v|\sin\alpha\big)$.}

\ex{1} Sean $\vec u=(2,\,-3)$, $\vec v=(-1,\,4)$ y $\vec w=(5,\,0)$. Calcula:\begin{tasks}(2)
  \task $\vec u+\vec v-\vec w$
  \task $3\vec u-2\vec v$
  \task $\tfrac12\vec w-\vec u+\vec v$
  \task $-2(\vec u-\vec v)+3\vec w$
\end{tasks}

\ex{2} Calcula $x$ e $y$ para que se cumpla $(x,\,3)+2(1,\,y)=(5,\,-1)$. Después, halla $k$ para que el vector $(k,\,4)$ tenga módulo 5.

\ex{2} Halla $m$ para que los vectores $\vec u=(m,\,-2)$ y $\vec v=(6,\,3)$ sean paralelos. Para ese valor, ¿qué número $k$ cumple $\vec u=k\vec v$?

\ex{2} Comprueba que $\vec u=(1,\,2)$ y $\vec v=(3,\,-1)$ forman una base y escribe $\vec w=(1,\,9)$ como combinación lineal de $\vec u$ y $\vec v$.


% =====================================================
\seccio{2}{Aplicaciones geométricas de los vectores}

\exf{1} Calcula los puntos medios de los lados del triángulo de vértices $A(-2,\,1)$, $B(4,\,3)$ y $C(2,\,-5)$.

\ex{1} El punto $M(1,\,-2)$ es el punto medio del segmento $AB$. Si $A(4,\,3)$, ¿cuáles son las coordenadas de $B$?

\ex{1} Halla los puntos que dividen el segmento de extremos $A(-3,\,4)$ y $B(9,\,-2)$ en tres partes iguales.

\ex{2} Divide el segmento de extremos $A(1,\,-1)$ y $B(11,\,4)$ en cinco partes iguales.

\ex{1} Estudia si los puntos están alineados:\begin{tasks}(2)
  \task $A(1,\,3)$, $B(3,\,7)$, $C(-2,\,-3)$
  \task $A(0,\,2)$, $B(3,\,4)$, $C(6,\,7)$
\end{tasks}Calcula $k$ para que $A(-1,\,2)$, $B(2,\,k)$ y $C(5,\,-4)$ estén alineados.

\ex{1} Calcula el simétrico del punto $P(2,\,-3)$ respecto del punto $Q(-1,\,4)$, y el simétrico de $R(5,\,1)$ respecto del origen.

\ex{2} Tres vértices consecutivos de un paralelogramo $ABCD$ son $A(-1,\,-2)$, $B(4,\,-1)$ y $C(5,\,3)$. Halla el vértice $D$ y el centro del paralelogramo.


% =====================================================
\seccio{3}{Producto escalar y ángulo de dos vectores}

\exf{1} Sean $\vec u=(3,\,-2)$ y $\vec v=(4,\,1)$. Calcula $\vec u\cdot\vec v$, $|\vec u|$, $|\vec v|$ y el ángulo que forman.

\ex{1} Calcula el ángulo que forman los vectores:\begin{tasks}(2)
  \task $(1,\,2)$ y $(-2,\,1)$
  \task $(1,\,1)$ y $(-1,\,0)$
  \task $(3,\,4)$ y $(-4,\,-3)$
  \task $(2,\,-1)$ y $(1,\,3)$
\end{tasks}

\ex{2} Si $|\vec u|=2$, $|\vec v|=5$ y el ángulo que forman es de $120^\circ$, calcula:\begin{tasks}(3)
  \task $\vec u\cdot\vec v$
  \task $(2\vec u+\vec v)\cdot(\vec u-\vec v)$
  \task $|\vec u+\vec v|$
\end{tasks}\recorda{$\vec u\cdot\vec u=|\vec u|^2$ y el producto escalar es distributivo.}

\ex{2} Halla $m$ para que $\vec u=(m,\,2)$ y $\vec v=(3,\,-6)$ sean perpendiculares. Halla también $k$ para que $(1,\,k)$ forme un ángulo de $60^\circ$ con $(1,\,0)$.

\ex{1} Halla los vectores unitarios perpendiculares a $\vec v=(6,\,-8)$.

\ex{2} Comprueba que el triángulo de vértices $A(1,\,1)$, $B(5,\,3)$ y $C(3,\,7)$ es rectángulo e isósceles. ¿Cuánto miden sus ángulos?

\ex{2} Calcula la proyección del vector $\vec u=(2,\,5)$ sobre la dirección de $\vec v=(3,\,-4)$. ¿Qué indica el signo?


% =====================================================
\seccio{4}{Ecuaciones de la recta}

\exf{1} Escribe todas las ecuaciones (vectorial, paramétricas, continua, punto-pendiente, explícita y general) de la recta que pasa por $P(2,\,-1)$ y tiene vector director $\vec u=(3,\,4)$.

\ex{1} Escribe todas las ecuaciones de la recta que pasa por $A(-1,\,3)$ y $B(3,\,-5)$.

\ex{1} Halla las ecuaciones punto-pendiente, explícita y general de la recta que pasa por $P(4,\,-2)$ y tiene pendiente $-\tfrac32$.

\ex{1} Dada la recta $r:2x-5y+10=0$, halla: la pendiente, los puntos de corte con los ejes, un vector director, un vector normal y unas ecuaciones paramétricas.

\ex{1} Comprueba si los puntos $P(7,\,1)$ y $Q(4,\,0)$ pertenecen a la recta $r:\left\{\begin{array}{l}x=1+3\lambda\\ y=-3+2\lambda\end{array}\right.$ Calcula $k$ para que $(k,\,5)$ sea de $r$.

\ex{1} Escribe la ecuación de la recta:\begin{tasks}(2)
  \task Paralela al eje $OX$ por $(3,\,-2)$.
  \task Paralela al eje $OY$ por $(3,\,-2)$.
  \task Que pasa por el origen y por $(3,\,-2)$.
  \task Bisectriz del 1.\textsuperscript{er} y 3.\textsuperscript{er} cuadrantes.
\end{tasks}

\ex{1} Halla la ecuación general de la recta que pasa por $P(1,\,5)$ y tiene vector normal $\vec n=(3,\,-1)$.

\ex{2} Halla la ecuación de la recta que pasa por $(2,\,3)$ y forma un ángulo de $135^\circ$ con el semieje positivo $OX$, y la de la recta que corta al eje $OY$ en $(0,\,-4)$ y forma un ángulo de $60^\circ$ con $OX$.


% =====================================================
\seccio{5}{Posiciones relativas de dos rectas}

\exf{1} Estudia la posición relativa de cada par de rectas y, si se cortan, halla el punto de corte:\begin{tasks}(1)
  \task $r:2x-3y+1=0$, \quad $s:4x+y-5=0$
  \task $r:x-2y+3=0$, \quad $s:-3x+6y-9=0$
  \task $r:3x+y-2=0$, \quad $s:6x+2y+1=0$
\end{tasks}

\ex{1} Estudia la posición relativa de $r:\left\{\begin{array}{l}x=1+2\lambda\\ y=-1+\lambda\end{array}\right.$ y $s:\left\{\begin{array}{l}x=3-4\mu\\ y=2\mu\end{array}\right.$

\ex{1} Estudia la posición relativa de $r:\left\{\begin{array}{l}x=2+\lambda\\ y=1-3\lambda\end{array}\right.$ y $s:3x+y-7=0$.

\ex{2} Halla $a$ para que las rectas $ax+3y-2=0$ y $2x-y+5=0$ sean: a) secantes; b) paralelas. ¿Pueden ser coincidentes?

\ex{2} Calcula $m$ y $n$ para que las rectas $mx+2y+n=0$ y $3x-y+4=0$ sean coincidentes.

\ex{2} Las rectas $r:x-y+1=0$, $s:x+y-5=0$ y $t:y=-1$ determinan un triángulo. Halla sus vértices.


% =====================================================
\seccio{6}{Paralelismo y perpendicularidad}

\exf{1} Halla las rectas paralela y perpendicular a $r:3x-2y+5=0$ que pasan por $P(2,\,-1)$.

\ex{1} Halla las rectas paralela y perpendicular a $r:\left\{\begin{array}{l}x=1-2\lambda\\ y=3+\lambda\end{array}\right.$ que pasan por $A(0,\,-2)$. Dalas en forma general.

\ex{1} Dada $r:y=-\tfrac13x+2$, halla las ecuaciones explícitas de las rectas paralela y perpendicular a $r$ que pasan por $(-3,\,4)$.

\ex{2} Calcula $k$ para que las rectas $kx-2y+1=0$ y $3x+(k+1)y-4=0$ sean perpendiculares. ¿Pueden ser paralelas?

\ex{2} En el triángulo de vértices $A(1,\,1)$, $B(7,\,3)$ y $C(3,\,7)$, halla la ecuación de la altura que sale de $C$ y el pie de esa altura.

\ex{2} Halla la mediatriz del segmento de extremos $A(-3,\,2)$ y $B(5,\,-4)$. Compruébala con la definición (puntos que equidistan de $A$ y de $B$).


% =====================================================
\seccio{7}{Ángulo de dos rectas}

\exf{1} Calcula el ángulo que forman las rectas $r:3x-y+2=0$ y $s:x+2y-1=0$.

\ex{1} Calcula el ángulo de $r:\left\{\begin{array}{l}x=1+2\lambda\\ y=3-\lambda\end{array}\right.$ y $s:\left\{\begin{array}{l}x=2-\lambda\\ y=4+3\lambda\end{array}\right.$

\ex{1} Calcula el ángulo de las rectas $y=2x-1$ e $y=-x+4$.

\ex{1} ¿Qué ángulo forma cada recta con el eje $OX$? \quad a) $3x-3y+1=0$ \quad b) $x+2y=0$

\ex{3} Halla las rectas que pasan por $P(1,\,2)$ y forman un ángulo de $45^\circ$ con la recta $r:x-3y+1=0$.


% =====================================================
\seccio{8}{Distancias}

\exf{1} Calcula las longitudes de los lados del triángulo de vértices $A(-1,\,2)$, $B(3,\,5)$ y $C(6,\,1)$ y clasifícalo. ¿Cuál es su perímetro?

\ex{1} Calcula la distancia del punto $P(3,\,-2)$ a cada una de las rectas:\begin{tasks}(2)
  \task $4x-3y+2=0$
  \task $y=2x+1$
  \task $x=-1$
  \task $\left\{\begin{array}{l}x=1+\lambda\\ y=2-\lambda\end{array}\right.$
\end{tasks}

\ex{1} Comprueba que las rectas $r:3x+4y-5=0$ y $s:6x+8y+15=0$ son paralelas y calcula la distancia entre ellas.

\ex{2} Halla los puntos del eje $OY$ que están a distancia 2 de la recta $3x-4y+1=0$.

\ex{2} Halla las rectas paralelas a $5x-12y+3=0$ que están a distancia 2 de ella.

\ex{2} Una circunferencia tiene el centro en $C(2,\,-1)$ y es tangente a la recta $3x+4y-12=0$. ¿Cuál es su radio?


% =====================================================
\seccio{9}{Problemas métricos}

\exf{2} Halla la proyección del punto $P(-3,\,4)$ sobre la recta $r:2x-y+5=0$ y el simétrico de $P$ respecto de $r$.

\ex{1} Calcula el simétrico del punto $P(4,\,1)$ respecto de la recta $y=x$ y respecto de la recta $x=2$.

\ex{3} Halla las bisectrices de los ángulos que forman las rectas $r:3x+4y-1=0$ y $s:5x-12y+2=0$. Comprueba que son perpendiculares.

\ex{2} Calcula el área del triángulo determinado por las rectas $r:x-y+1=0$, $s:x+y-5=0$ y $t:y=-1$ (ejercicio 35).

\ex{2} En el triángulo $A(-2,\,0)$, $B(4,\,2)$, $C(1,\,6)$, calcula:\begin{tasks}(2)
  \task La mediana que sale de $A$.
  \task La altura que sale de $A$.
  \task El baricentro.
  \task El área.
\end{tasks}

\ex{2} Halla el circuncentro y el radio de la circunferencia circunscrita al triángulo $A(0,\,0)$, $B(6,\,0)$, $C(0,\,4)$.

\ex{2} Halla el ortocentro del triángulo de vértices $A(-3,\,0)$, $B(3,\,0)$ y $C(1,\,4)$.

\ex{2} Los puntos $A(1,\,1)$ y $C(5,\,5)$ son vértices opuestos de un cuadrado. Halla los otros dos vértices y el área del cuadrado.

\ex{3} Halla el punto de la recta $y=x$ que está a la misma distancia de $A(1,\,0)$ que de $B(5,\,2)$.

\vfill
\begin{center}\color{muted}\small\faCheckCircle[regular]\ \ ¡Ánimo! Comprueba tus resultados con las herramientas interactivas de los apuntes de la unidad.\end{center}

\end{document}
